\chapter{How D'Artagnan drew a Country-Seat from a Deal Box}

\lettrine{T}{he} king's words regarding the wounded pride of Monk had inspired D'Artagnan with no small portion of apprehension. The lieutenant had had, all his life, the great art of choosing his enemies; and when he had found them implacable and invincible, it was when he had not been able, under any pretence, to make them otherwise. But points of view change greatly in the course of a life. It is a magic lantern, of which the eye of man every year changes the aspects. It results that from the last day of a year on which we saw white, to the first day of the year on which we shall see black, there is the interval of but a single night.

Now, D'Artagnan, when he left Calais with his ten scamps, would have hesitated as little in attacking a Goliath, a Nebuchadnezzar, or a Holofernes, as he would in crossing swords with a recruit or caviling with a land-lady. Then he resembled the sparrow-hawk, which, when fasting, will attack a ram. Hunger is blind. But D'Artagnan satisfied—D'Artagnan rich—D'Artagnan a conqueror—D'Artagnan proud of so difficult a triumph—D'Artagnan had too much to lose not to reckon, figure by figure, with probable misfortune.

His thoughts were employed, therefore, all the way on the road from his presentation, with one thing, and that was, how he should conciliate a man like Monk, a man whom Charles himself, king as he was, conciliated with difficulty; for, scarcely established, the protected might again stand in need of the protector, and would, consequently, not refuse him, such being the case, the petty satisfaction of transporting M.~d'Artagnan, or of confining him in one of the Middlesex prisons, or drowning him a little on his passage from Dover to Boulogne. Such sorts of satisfaction kings are accustomed to render to viceroys without disagreeable consequences.

It would not be at all necessary for the king to be active in that contrepartie of the play in which Monk should take his revenge. The part of the king would be confined to simply pardoning the viceroy of Ireland all he should undertake against D'Artagnan. Nothing more was necessary to place the conscience of the Duke of Albemarle at rest than a \binderylatin{te absolvo} said with a laugh, or the scrawl of <Charles the King,> traced at the foot of a parchment; and with these two words pronounced, and these two words written, poor D'Artagnan was forever crushed beneath the ruins of his imagination.

And then, a thing sufficiently disquieting for a man with such foresight as our musketeer, he found himself alone; and even the friendship of Athos could not restore his confidence. Certainly if the affair had only concerned a free distribution of sword-thrusts, the musketeer would have counted upon his companion; but in delicate dealings with a king, when the perhaps of an unlucky chance should arise in justification of Monk or of Charles of England, D'Artagnan knew Athos well enough to be sure he would give the best possible colouring to the loyalty of the survivor, and would content himself with shedding floods of tears on the tomb of the dead, supposing the dead to be his friend, and afterwards composing his epitaph in the most pompous superlatives.

<Decidedly,> thought the Gascon; and this thought was the result of the reflections which he had just whispered to himself and which we have repeated aloud—<decidedly, I must be reconciled with M.~Monk, and acquire proof of his perfect indifference for the past. If, and God forbid it should be so! he is still sulky and reserved in the expression of this sentiment, I shall give my money to Athos to take away with him, and remain in England just long enough to unmask him, then, as I have a quick eye and a light foot, I shall notice the first hostile sign; to decamp or conceal myself at the residence of my lord Buckingham, who seems a good sort of devil at the bottom, and to whom, in return for his hospitality, I shall relate all that history of the diamonds, which can now compromise nobody but an old queen, who need not be ashamed, after being the wife of a miserly creature like Mazarin, of having formerly been the mistress of a handsome nobleman like Buckingham. \swear{Mordioux!} that is the thing, and this Monk shall not get the better of me. Eh? and besides I have an idea!>

We know that, in general, D'Artagnan was not wanting in ideas; and during this soliloquy, D'Artagnan buttoned his vest up to the chin, and nothing excited his imagination like this preparation for a combat of any kind, called accinction by the Romans. He was quite heated when he reached the mansion of the Duke of Albemarle. He was introduced to the viceroy with a promptitude which proved that he was considered as one of the household. Monk was in his business-closet.

<My lord,> said D'Artagnan, with that expression of frankness which the Gascon knew so well how to assume, <my lord, I have come to ask your grace's advice!>

Monk, as closely buttoned up morally as his antagonist was physically, replied: <Ask, my friend;> and his countenance presented an expression not less open than that of D'Artagnan.

<My lord, in the first place, promise me secrecy and indulgence.>

<I promise you all you wish. What is the matter? Speak!>

<It is, my lord, that I am not quite pleased with the king.>

<Indeed! And on what account, my dear lieutenant?>

<Because his majesty gives way sometimes to jests very compromising for his servants; and jesting, my lord, is a weapon that seriously wounds men of the sword, as we are.>

Monk did all in his power not to betray his thought, but D'Artagnan watched him with too close attention not to detect an almost imperceptible flush upon his face. <Well, now, for my part,> said he, with the most natural air possible, <I am not an enemy of jesting, my dear Monsieur d'Artagnan; my soldiers will tell you that even many times in camp, I listened very indifferently, and with a certain pleasure, to the satirical songs which the army of Lambert passed into mine, and which, certainly, would have caused the ears of a general more susceptible than I am to tingle.>

<Oh, my lord,> said D'Artagnan, <I know you are a complete man; I know you have been, for a long time, placed above human miseries; but there are jests and jests of a certain kind, which have the power of irritating me beyond expression.>

<May I inquire what kind, my friend?>

<Such as are directed against my friends or against people I respect, my lord!>

Monk made a slight movement, which D'Artagnan perceived. <Eh! and in what,> asked Monk, <in what can the stroke of a pin which scratches another tickle your skin? Answer me that.>

<My lord, I can explain it to you in a single sentence; it concerns you.>

Monk advanced a single step towards D'Artagnan. <Concerns me?> said he.

<Yes, and this is what I cannot explain; but that arises, perhaps, from my want of knowledge of his character. How can the king have the heart to jest about a man who has rendered him so many and such great services? How can one understand that he should amuse himself in setting by the ears a lion like you with a gnat like me?>

<I cannot conceive that in any way,> said Monk.

<But so it is. The king, who owed me a reward, might have rewarded me as a soldier, without contriving that history of the ransom, which affects you, my lord.>

<No,> said Monk, laughing: <it does not affect me in any way, I can assure you.>

<Not as regards me, I can understand; you know me, my lord, I am so discreet that the grave would appear a babbler compared to me; but—do you understand, my lord?>

<No,> replied Monk, with persistent obstinacy.

<If another knew the secret which I know\longdash>

<What secret?>

<Eh! my lord, why, that unfortunate secret of Newcastle.>

<Oh! the million of the Comte de la Fère?>

<No, my lord, no; the enterprise made upon your grace's person.>

<It was well played, chevalier, that is all, and no more is to be said about it: you are a soldier, both brave and cunning, which proves that you unite the qualities of Fabius and Hannibal. You employed your means, force and cunning: there is nothing to be said against that: I ought to have been on guard.>

<Ah! yes; I know, my lord, and I expected nothing less from your partiality; so that if it were only the abduction in itself, \swear{Mordioux!} that would be nothing; but there are\longdash>

<What?>

<The circumstances of that abduction.>

<What circumstances?>

<Oh! you know very well what I mean, my lord.>

<No, curse me if I do.>

<There is—in truth, it is difficult to speak it.>

<There is?>

<Well, there is that devil of a box!>

Monk coloured visibly. <Well, I have forgotten it.>

<Deal box,> continued D'Artagnan, <with holes for the nose and mouth. In truth, my lord, all the rest was well; but the box, the box! that was really a coarse joke.> Monk fidgeted about in his chair. <And, notwithstanding my having done that,> resumed D'Artagnan, <I, a soldier of fortune, it was quite simple, because by the side of that action, a little inconsiderate I admit, which I committed, but which the gravity of the case may excuse, I am circumspect and reserved.>

<Oh!> said Monk, <believe me, I know you well, Monsieur d'Artagnan, and I appreciate you.>

D'Artagnan never took his eyes off Monk; studying all which passed in the mind of the general, as he prosecuted his idea. <But it does not concern me,> resumed he.

<Well, then, who does it concern?> said Monk, who began to grow a little impatient.

<It relates to the king, who will never restrain his tongue.>

<Well! and suppose he should say all he knows?> said Monk, with a degree of hesitation.

<My lord,> replied D'Artagnan, <do not dissemble, I implore you, with a man who speaks so frankly as I do. You have a right to feel your susceptibility excited, however benignant it may be. What, the devil! it is not the place for a man like you, a man who plays with crowns and sceptres as a Bohemian plays with his balls; it is not the place of a serious man, I said, to be shut up in a box like some freak of natural history; for you must understand it would make all your enemies ready to burst with laughter, and you are so great, so noble, so generous, that you must have many enemies. This secret is enough to set half the human race laughing, if you were represented in that box. It is not decent to have the second personage in the kingdom laughed at.>

Monk was quite out of countenance at the idea of seeing himself represented in this box. Ridicule, as D'Artagnan had judiciously foreseen, acted upon him in a manner which neither the chances of war, the aspirations of ambition, nor the fear of death had been able to do.

<Good,> thought the Gascon, <he is frightened: I am safe.>

<Oh! as to the king,> said Monk, <fear nothing, my dear Monsieur d'Artagnan; the king will not jest with Monk, I assure you!>

The momentary flash of his eye was noticed by D'Artagnan. Monk lowered his tone immediately: <The king,> continued he, <is of too noble a nature, the king's heart is too high to allow him to wish ill to those who do him good.>

<Oh! certainly,> cried D'Artagnan. <I am entirely of your grace's opinion with regard to his heart, but not as to his head—it is good, but it is trifling.>

<The king will not trifle with Monk, be assured.>

<Then you are quite at ease, my lord?>

<On that side, at least! yes, perfectly!>

<Oh! I understand you; you are at ease as far as the king is concerned?>

<I have told you I was.>

<But you are not so much so on my account?>

<I thought I had told you that I had faith in your loyalty and discretion.>

<No doubt, no doubt, but you must remember one thing\longdash>

<What is that?>

<That I was not alone, that I had companions; and what companions!>

<Oh! yes, I know them.>

<And, unfortunately, my lord, they know you, too!>

<Well?>

<Well; they are yonder, at Boulogne, waiting for me.>

<And you fear\longdash>

<Yes, I fear that in my absence—\swear{Parbleu!} If I were near them, I could answer for their silence.>

<Was I not right in saying that the danger, if there was any danger, would not come from his majesty, however disposed he may be to jest, but from your companions, as you say? To be laughed at by a king may be tolerable, but by the horse-boys and scamps of the army! Damn it!>

<Yes, I understand, that would be unbearable; that is why, my lord, I came to say,—do you not think it would be better for me to set out for France as soon as possible?>

<Certainly, if you think your presence\longdash>

<Would impose silence upon those scoundrels? Oh! I am sure of that, my lord.>

<Your presence will not prevent the report from spreading, if the tale has already transpired.>

<Oh! it has not transpired, my lord, I will wager. At all events, be assured that I am determined upon one thing.>

<What is that?>

<To blow out the brains of the first who shall have propagated that report, and of the first who has heard it. After which I shall return to England to seek an asylum, and perhaps employment with your grace.>

<Oh, come back! come back!>

<Unfortunately, my lord, I am acquainted with nobody here but your grace, and if I should no longer find you, or if you should have forgotten me in your greatness?>

<Listen to me, Monsieur d'Artagnan,> replied Monk; <you are a superior man, full of intelligence and courage; you deserve all the good fortune this world can bring you; come with me into Scotland, and, I swear to you, I shall arrange for you a fate which all may envy.>

<Oh! my lord, that is impossible. At present I have a sacred duty to perform; I have to watch over your glory, I have to prevent a low jester from tarnishing in the eyes of our contemporaries—who knows? in the eyes of posterity—the splendour of your name.>

<Of posterity, Monsieur d'Artagnan?>

<Doubtless. It is necessary, as regards posterity, that all the details of that history should remain a mystery; for, admit that this unfortunate history of the deal box should spread, and it should be asserted that you had not re-established the king loyally, and of your own free will, but in consequence of a compromise entered into at Scheveningen between you two. It would be vain for me to declare how the thing came about, for though I know I should not be believed, it would be said that I had received my part of the cake, and was eating it.>

Monk knitted his brow.—<Glory, honour, probity!> said he, <you are but empty words.>

<Mist!> replied D'Artagnan; <nothing but mist, through which nobody can see clearly.>

<Well, then, go to France, my dear Monsieur d'Artagnan,> said Monk; <go, and to render England more attractive and agreeable to you, accept a remembrance of me.>

<What now?> thought D'Artagnan.

<I have on the banks of the Clyde,> continued Monk, <a little house in a grove, cottage as it is called here. To this house are attached a hundred acres of land. Accept it as a souvenir.>

<Oh, my lord!—>

<Faith! you will be there in your own home, and that will be the place of refuge you spoke of just now.>

<For me to be obliged to your lordship to such an extent! Really, your grace, I am ashamed.>

<Not at all, not at all, monsieur,> replied Monk, with an arch smile; <it is I who shall be obliged to you. And,> pressing the hand of the musketeer, <I shall go and draw up the deed of gift,>—and he left the room.

D'Artagnan looked at him as he went out with something of a pensive and even an agitated air.

<After all,> said he, <he is a brave man. It is only a sad reflection that it is from fear of me, and not affection that he acts thus. Well, I shall endeavour that affection may follow.> Then, after an instant's deeper reflection,—<Bah!> said he, <to what purpose? He is an Englishman.> And he in turn went out, a little confused after the combat.

<So,> said he, <I am a land-owner! But how the devil am I to share the cottage with Planchet? Unless I give him the land, and I take the château, or the he takes the house and I—nonsense! M.~Monk will never allow me to share a house he has inhabited, with a grocer. He is too proud for that. Besides, why should I say anything about it to him? It was not with the money of the company I have acquired that property, it was with my mother-wit alone; it is all mine, then. So, now I will go and find Athos.> And he directed his steps towards the dwelling of the Comte de la Fère. 